% ThunderServe 4-minute summary. Input from 2026-10-28-Meeting-Main.tex.
% Figures are copies of 2026-10-21-meeting/ThunderServe/. See SOURCE.txt.
% Facts checked against 2026-10-21-meeting/MLSys25_Thunderserve.pdf.

\begin{frame}[t]{Why Heterogeneous GPU Serving Needs Phase-Aware Scheduling}
\begin{columns}[T,onlytextwidth]
\begin{column}{.50\textwidth}
\mainfig{ThunderServe/ThunderServe_Fig01.png}{3.85cm}
\end{column}
\begin{column}{.48\textwidth}
{\footnotesize\textbf{Observation}}
\begin{itemize}\setlength{\itemsep}{.04cm}\footnotesize
\item Prefill: compute-intensive.
\item Decode: memory-bandwidth-intensive.
\end{itemize}
\vspace{.06cm}
{\scriptsize A40: 149.7 FP16 TFLOPS. 3090 Ti: 1008 GB/s. Table~1.}
\vspace{.12cm}
{\footnotesize\textbf{Challenge}}
\begin{itemize}\setlength{\itemsep}{.04cm}\footnotesize
\item Heterogeneous GPU resources.
\item Limited and heterogeneous network bandwidth.
\item KV cache transfer overhead.
\end{itemize}
\vspace{.08cm}
{\footnotesize\textbf{Research question}\par}
{\footnotesize How to jointly optimize GPU allocation, phase placement, and communication?\par}
\end{column}
\end{columns}
\vspace{.08cm}
\takeaway{The better GPU can differ by phase. KV transfer makes a FLOPS-only or bandwidth-only split insufficient.}
\src{ThunderServe Sec.~1, Fig.~1, Table~1. Prices in the figure are per request.}
\end{frame}

\begin{frame}[t]{Joint Optimization of GPU Placement, Parallelism, and Routing}
\begin{columns}[T,onlytextwidth]
\begin{column}{.60\textwidth}
\mainfig{ThunderServe/ThunderServe_Fig03.png}{4.70cm}
\end{column}
\begin{column}{.38\textwidth}
\smalllabel{Upper-level --- Tabu search}\par
\begin{itemize}\setlength{\itemsep}{.04cm}\footnotesize
\item GPU group construction.
\item Prefill / decode phase designation.
\end{itemize}
\vspace{.08cm}
\smalllabel{Lower-level --- configuration}\par
\begin{itemize}\setlength{\itemsep}{.04cm}\footnotesize
\item TP / PP configuration.
\item Prefill / decode request routing.
\item Network-aware KV cache transfer.
\end{itemize}
\vspace{.08cm}
{\scriptsize Additional mechanisms: KV cache compression and lightweight rescheduling.}
\end{column}
\end{columns}
\vspace{.04cm}
\takeaway{Upper level assigns groups and phases. Lower level sets parallelism and routing. Together they score one deployment plan.}
\src{ThunderServe Sec.~3.2--3.4, Fig.~3.}
\end{frame}

\begin{frame}[t]{Does Phase-Aware Scheduling Improve Serving Efficiency?}
{\scriptsize\textbf{Question:} Does workload-aware phase placement raise end-to-end throughput?\par}
\vspace{.06cm}
\begin{columns}[T,onlytextwidth]
\begin{column}{.46\textwidth}
\smalllabel{Experiment 1 --- Table 3}\par
{\footnotesize\textbf{Workload-aware deployment}\par}
\vspace{.04cm}
{\footnotesize Coding: \textbf{8 prefill / 4 decode}\par}
{\footnotesize Conversation: \textbf{4 prefill / 8 decode}\par}
\vspace{.06cm}
{\scriptsize LLaMA-30B, same 32-GPU cloud pool. Coding has longer prompts; conversation has longer outputs.\par}
\vspace{.10cm}
\smalllabel{Experiment 2 --- Figure 9}\par
\vspace{.04cm}
{\scriptsize\renewcommand{\arraystretch}{1.08}
\begin{tabular}{@{}lcc@{}}\toprule
 & \textbf{Coding} & \textbf{Conv.}\\\midrule
vs.\ HexGen & $1.5\times$ & $1.3\times$\\
vs.\ DistServe & $1.5\times$ & $2.1\times$\\\bottomrule
\end{tabular}\par}
\vspace{.05cm}
{\scriptsize HexGen: heterogeneous cloud GPUs.\par
DistServe: homogeneous 8$\times$A100, similar price. Not the same machines.\par}
\end{column}
\begin{column}{.52\textwidth}
\mainfig{ThunderServe/ThunderServe_Fig09.png}{3.35cm}
\vspace{.06cm}
{\scriptsize vLLM, the fourth group, is also in-house. HexGen and DistServe are different hardware setups.\par}
\end{column}
\end{columns}
\vspace{.08cm}
\takeaway{Workload-aware placement improves overall serving throughput under heterogeneous resource constraints.}
\src{ThunderServe Sec.~5.1--5.3, Fig.~9, Table~3. Relative throughput is ThunderServe over each baseline.}
\end{frame}

\begin{frame}[t]{Phase-Aware Deployment Improves Cost Efficiency}
\vspace{.08cm}
\begin{columns}[T,onlytextwidth]
\begin{column}{.32\textwidth}
\stepcard{1. JOINT OPTIMIZATION}{One deployment plan}{GPU grouping, phase designation, TP/PP, and request routing.}
\end{column}
\begin{column}{.32\textwidth}
\stepcard{2. PERFORMANCE}{1.7$\times$ average, up to 2.1$\times$}{Overall throughput claim in the paper. Not one bar of Figure~9.}
\end{column}
\begin{column}{.32\textwidth}
\stepcard{3. DYNAMIC ADAPTATION}{Full 157 s, lightweight 13 s}{Table~4 overall time. The gap includes model reload, not search time alone.}
\end{column}
\end{columns}
\vspace{.16cm}
{\scriptsize Table~4 breakdown. Full: rescheduling 54($\pm$5)\,s $+$ reload 103($\pm$10)\,s $=$ 157($\pm$13)\,s.
Lightweight: rescheduling 13($\pm$2)\,s $+$ reload 0\,s $=$ 13($\pm$2)\,s.\par}
\vspace{.12cm}
\takeaway{ThunderServe demonstrates that heterogeneous GPU serving benefits from jointly optimizing inference phases, hardware allocation, and communication.}
\vspace{.08cm}
{\footnotesize ThunderServe focuses on phase-level resource allocation; finer-grained workload partitioning is another possible direction.\par}
\src{Abstract and Sec.~6 for the 1.7$\times$/2.1$\times$ claim. Sec.~5.4 and Table~4 for rescheduling time.}
\end{frame}
